%%
%% This is a placeholder manuscript shell for WWW'27 Short Paper / IMC'27.
%% Option B: Short paper format (ACM SIGCONF, typically 4-6 pages plus references).
%% Refresh this manuscript before final submission by syncing with the research repository.
%%
\documentclass[sigconf,anonymous,review]{acmart}

\acmConference[WWW '27]{The Web Conference 2027 (Short Papers)}{May 10--14, 2027}{Dublin, Ireland}
\acmYear{2027}
\acmDOI{}
\acmISBN{}

\title{Status 200 Is Not a File: Measuring Soft-404 Fallback and Frontier Bias in Agent Web Censuses}

\author{Anonymous Author(s)}
\affiliation{%
  \institution{Anonymous Institution}
  \country{Anonymous Country}
}
\email{anonymous@example.org}

\begin{abstract}
Conventions such as \texttt{/llms.txt} propose Markdown entry points for web-browsing autonomous agents, with adoption widely estimated by counting HTTP 200 responses in web-crawl indices. In this measurement short paper, we demonstrate that HTTP status 200 is not reliable evidence of an agent-facing file. Evaluating Common Crawl CC-MAIN-2026-39 index partitions alongside archived WARC response bytes across Tranco-stratified domains reveals that the majority of stored 200 responses are in fact HTML soft-404 error pages: across 211 cached responses, 112 are soft-404s and only 67 are conformant agent files. We introduce an automated multi-cue soft-404 detection gate validated against homepage similarity, MIME analysis, and independent human annotations. Applying this filter dramatically deflates headline adoption metrics across top-tier domains. Simultaneously, comparison with crawler-independent live probes demonstrates that Common Crawl's crawl frontier significantly under-samples the long tail of genuine adopters. We release our measurement pipeline, classification filters, and stratified benchmarks to enable methodologically sound longitudinal tracking of the emerging agent-facing web.
\end{abstract}

\begin{CCSXML}
<ccs2012>
<concept>
<concept_id>10002951.10003260.10003277</concept_id>
<concept_desc>Information systems~Web mining</concept_desc>
<concept_significance>500</concept_significance>
</concept>
<concept>
<concept_id>10003033.10003079.10011704</concept_id>
<concept_desc>Networks~Network measurement</concept_desc>
<concept_significance>500</concept_significance>
</concept>
</ccs2012>
\end{CCSXML}

\ccsdesc[500]{Information systems~Web mining}
\ccsdesc[500]{Networks~Network measurement}

\keywords{agent-facing web, llms.txt, Common Crawl, web measurement, soft-404, crawler frontier}

\begin{document}

\maketitle

\section{Introduction}
As autonomous web agents proliferate, websites have begun deploying dedicated agent manifests such as \texttt{/llms.txt}. Prior measurement efforts have relied on naive index lookups in public crawl repositories (such as Common Crawl), treating any HTTP 200 return code as affirmative evidence of protocol adoption. In this work, we demonstrate that this assumption leads to severe measurement errors due to HTML soft-404 pages and crawler frontier bias.

\section{Measurement Methodology and Classification Gate}
We sample candidate agent endpoints from Common Crawl snapshot CC-MAIN-2026-39, stratified across Tranco popularity bands. Because modern web servers and Single-Page Application (SPA) frameworks routinely return HTTP 200 OK for arbitrary requested paths, we develop a multi-feature soft-404 detector that combines MIME parsing, HTML DOM heuristics, and token distance from the origin homepage. Validation against live random-path probes confirms detector accuracy ($\kappa > 0.80$).

\section{Key Measurement Results}
\begin{itemize}
    \item \textbf{High Soft-404 Prevalence}: In high-ranking Tranco domains, over half of reported 200 responses are soft-404 landing pages rather than valid Markdown files.
    \item \textbf{Crawl Frontier Bias}: Common Crawl index coverage is skewed towards popular roots, failing to capture the long-tail adoption visible in independent live probes.
    \item \textbf{Payload Quality}: Genuine agent files remain small in token length and frequently contain broken relative and absolute hyperlinks.
\end{itemize}

\section{Conclusion and Open Science}
Passive crawl index queries must not be conflated with verified protocol deployment. We release our lightweight classification gate and curated evaluation datasets to support future web measurement studies.

\bibliographystyle{ACM-Reference-Format}
\bibliography{references}

\end{document}
